\documentclass[11pt]{article}

\usepackage[utf8]{inputenc}
\usepackage[T1]{fontenc}
\usepackage{lmodern}
\usepackage[margin=1in]{geometry}
\usepackage{amsmath}
\usepackage{amssymb}
\usepackage{graphicx}
\usepackage{booktabs}
\usepackage{xcolor}
\usepackage{listings}
\usepackage[numbers,sort&compress]{natbib}
\usepackage[hidelinks]{hyperref}

% Rust is not a built-in listings language; define a minimal one.
\lstdefinelanguage{Rust}{
  keywords={let, mut, fn, use, pub, struct, impl, trait, for, in, if, else,
    match, Some, None, true, false, self, Self, as, return, where, dyn, mod,
    const, move},
  morecomment=[l]{//},
  morecomment=[s]{/*}{*/},
  morestring=[b]",
  sensitive=true,
}

% Minimal, robust code style (Python / Rust).
\definecolor{codebg}{gray}{0.96}
\lstdefinestyle{code}{
  basicstyle=\ttfamily\footnotesize,
  backgroundcolor=\color{codebg},
  breaklines=true,
  frame=single,
  framerule=0pt,
  columns=fullflexible,
  keepspaces=true,
  showstringspaces=false,
  xleftmargin=4pt,
  xrightmargin=4pt,
  aboveskip=8pt,
  belowskip=8pt,
}
\lstset{style=code}

\title{\textbf{turboswarm}: A general-purpose, extensible Particle Swarm
Optimization library with a Rust core}

\author{Jose L. Salmeron\thanks{CUNEF Universidad, Madrid, Spain. ORCID:
\href{https://orcid.org/0000-0001-7811-3716}{0000-0001-7811-3716}.}}

\date{June 2026}

\begin{document}

\maketitle

\begin{abstract}
\noindent
\texttt{turboswarm} is a general-purpose, open-source Particle Swarm
Optimization (PSO) library with a compute core written in Rust and an ergonomic
Python API. It targets real, integer, mixed, binary and interval-valued
(``grey'') decision variables, and exposes established algorithmic variants,
swarm topologies, constraint handling, multi-objective optimization, and
visualization behind a single, uniform interface. The optimization loop is
fully decoupled from any concrete variant through three traits
(\texttt{SearchSpace}, \texttt{Velocity}, \texttt{Topology}), so adding a new
variant is a self-contained trait implementation rather than a fork of the core.
The same Rust core is usable directly from Rust through zero-cost generics and
from Python through PyO3 and \texttt{maturin}, where native objectives run
without holding the Global Interpreter Lock. Every experiment is seedable and
therefore reproducible. On a five-function benchmark across three dimensions,
the compiled core is the fastest among the libraries compared, with a measured
speedup of roughly $2$--$3.5\times$ over the fastest pure-Python competitor at
comparable solution quality. The library is distributed on PyPI
(\texttt{turboswarm}) and crates.io (\texttt{turboswarm-core}), documented, and
archived on Zenodo (DOI \texttt{10.5281/zenodo.20832446}).
\end{abstract}

\section{Introduction}

Particle Swarm Optimization \citep{kennedy1995particle} is a population-based,
gradient-free optimization method in which a swarm of candidate solutions
(``particles'') moves through the search space guided by each particle's own
best position and the best position found in its neighbourhood. It is one of the
most widely used metaheuristics for continuous and combinatorial optimization,
with applications across engineering, machine learning and the sciences.

Researchers and students need a PSO implementation that is fast, reproducible,
broad in scope, and easy to extend with new variants for comparative studies.
Existing Python libraries cover parts of this need, but they are implemented
predominantly in pure Python, which ties the cost of the optimization loop to
the interpreter and makes large or repeated experiments slow. \texttt{turboswarm}
was written to provide a compiled, reproducible and extensible PSO toolkit that
is nonetheless as easy to call as a pure-Python package, suitable for research
experiments, algorithm benchmarking and teaching.

\section{State of the field}

Several mature open-source PSO and metaheuristics packages exist, including
\texttt{pyswarms} \citep{miranda2018pyswarms}, \texttt{pymoo}
\citep{blank2020pymoo}, \texttt{DEAP} \citep{fortin2012deap} and
\texttt{nevergrad} \citep{rapin2018nevergrad}. These are valuable and widely
adopted, but they are implemented mainly in pure Python, so the swarm loop
itself runs in the interpreter. \texttt{turboswarm} differs in three ways that
motivated building a new library rather than contributing to an existing one.
First, its optimization loop and native objectives are compiled (Rust) and run
without the Global Interpreter Lock (GIL), giving a measured speed advantage
(Section~\ref{sec:perf}). Second, it treats extensibility as a first-class
design goal: a new variant is a self-contained trait implementation, not a fork
of the core (Section~\ref{sec:design}). Third, it provides a grey-number
(interval-valued) search space, in which each decision variable is known only to
lie within an interval --- a capability we are not aware of in the libraries
above. Together these make \texttt{turboswarm} complementary to, rather than a
re-implementation of, the existing ecosystem.

\section{Software design}
\label{sec:design}

The optimization loop is fully decoupled from any concrete variant through three
traits.

\begin{itemize}
  \item \texttt{SearchSpace} defines the domain. The only difference between
  real and integer variables lives in its \texttt{decode} step, so the swarm
  always operates on a continuous internal representation
  ($\mathbb{R}^n$) and discretization happens only at evaluation time.
  \item \texttt{Velocity} defines the update rule, so a single PSO variant ---
  inertia \citep{shi1998modified}, constriction \citep{clerc2002particle} or the
  Fully Informed Particle Swarm \citep{mendes2004fully} --- is exactly one
  implementation of this trait.
  \item \texttt{Topology} defines the social structure of the swarm: global
  (gbest), ring (lbest), Von Neumann or random.
\end{itemize}

Because the loop is generic over these traits, adding a variant requires no
change to the core and is accompanied by a convergence test against a known
optimum. For the Python boundary the same traits are implemented for boxed trait
objects, so variants can be selected at runtime by name without duplicating the
loop. This design confines the real/integer/mixed/binary/grey distinction to the
search space and makes the library a convenient base for comparing variants
under identical conditions.

\section{Features}

\begin{itemize}
  \item \textbf{Velocity variants:} inertia \citep{shi1998modified},
  constriction \citep{clerc2002particle} and Fully Informed PSO
  \citep{mendes2004fully}.
  \item \textbf{Topologies:} global (gbest), ring (lbest), Von Neumann and
  random.
  \item \textbf{Search spaces:} continuous, integer (with a discretization
  strategy), binary, mixed (per-dimension type) and grey (interval-valued
  variables, optimized over their centre and spread).
  \item \textbf{Multi-objective (MOPSO):} an external Pareto archive with a
  turbulence operator, crowding-distance or Coello's adaptive grid
  \citep{coello2004handling} for diversity, and a hypervolume quality indicator.
  \item \textbf{Constraints:} inequality and equality constraints via a penalty,
  plus an optional repair operator that maps candidates back to the feasible set.
  \item \textbf{Run control:} stopping on a target value, evaluation budget,
  wall-clock budget or stagnation, with a per-iteration callback and a reported
  stop reason.
  \item \textbf{Boundary handling:} clamp, reflect, wrap or reinitialize.
  \item \textbf{Performance:} velocity clamping, \texttt{rayon}-based parallel
  evaluation and a NumPy-vectorized batch path.
  \item \textbf{Hyperparameter sensitivity:} a Cartesian-product sweep over
  hyperparameter values, aggregated over seeds.
  \item \textbf{Benchmarks:} sphere, rastrigin, rosenbrock, ackley, griewank,
  schwefel and the CEC-family functions (bent cigar, discus, elliptic, zakharov,
  levy, expanded Schaffer), each with metadata (recommended bound, known
  optimum).
  \item \textbf{Visualization (Python):} convergence curves, variant comparison,
  2D and 3D swarm animation, Pareto-front and sensitivity plots.
  \item \textbf{Integrations:} a SciPy \texttt{minimize} drop-in, a
  scikit-learn \texttt{PSOSearchCV}, an Optuna sampler, pandas export and
  parallel-objective backends.
\end{itemize}

\section{Usage}

From Python, an objective can be a callable or the name of a native benchmark;
bounds are given per dimension or as a single pair broadcast over \texttt{dim}.

\begin{lstlisting}[language=Python]
import turboswarm as pso

# Native benchmark, runs in Rust without the GIL
r = pso.minimize("rastrigin", bounds=(-5.12, 5.12), dim=2, seed=42)

# A user objective, a chosen variant and topology
r = pso.minimize(lambda x: sum(xi*xi for xi in x),
                 bounds=[(-5, 5)] * 3, velocity="fips", topology="ring", seed=1)
print(r.best_position, r.best_value)
\end{lstlisting}

The same core is callable from Rust through zero-cost generics:

\begin{lstlisting}[language=Rust]
use turboswarm_core::prelude::*;
use turboswarm_core::benchmarks::rastrigin;

let space = ContinuousSpace::uniform(2, -5.12, 5.12);
let velocity = InertiaVelocity::new(0.9, 1.49445, 1.49445).with_decay(0.4);
let params = PsoParams { seed: Some(42), ..Default::default() };
let result = Pso::new(space, velocity, GlobalBest::new(), params).minimize(rastrigin);
\end{lstlisting}

\section{Performance}
\label{sec:perf}

To illustrate the benefit of the compiled core we compare \texttt{turboswarm}
against four widely used baselines --- \texttt{pyswarms}
\citep{miranda2018pyswarms}, \texttt{pyswarm}, \texttt{pymoo}
\citep{blank2020pymoo} and DEAP \citep{fortin2012deap} --- on five standard
functions (sphere, rastrigin, ackley, rosenbrock, griewank) across three
dimensions (2, 10, 30), using an identical swarm configuration for every library
(40 particles, 200 iterations, $w = 0.729$, $c_1 = c_2 = 1.49445$; median over
five seeds after a warm-up). DEAP, a general-purpose evolutionary-computation
toolkit, provides a PSO \emph{recipe} rather than a packaged solver, which we
adapted to the same inertia rule and configuration. We report both
\texttt{turboswarm}'s native objective path (objective evaluated in Rust) and its
Python-callable path (objective evaluated in Python).

Table~\ref{tab:time} reports the median wall-clock time over the five functions
at each dimension, and Table~\ref{tab:speedup} the resulting speedup of the
native path over every baseline. Through its native path \texttt{turboswarm} is
the fastest library at every dimension: about $3.2\times$ faster than
\texttt{pyswarms} (the fastest competitor) at two dimensions, narrowing to
$2.1\times$ at thirty as \texttt{pyswarms}' vectorized NumPy objective amortizes
more of its per-iteration overhead. Against the interpreter-bound,
non-vectorized baselines the gap is larger and persistent: a median of about
$12\times$ over DEAP, $7$--$14\times$ over \texttt{pyswarm}, and $24$--$70\times$
over \texttt{pymoo}. Even \texttt{turboswarm}'s own Python-callable path is
competitive with the fastest baseline at low dimension. Solution quality was
comparable across all libraries --- the median best objective values fell within
the same order of magnitude (for example $0.08$--$0.11$ at ten dimensions) --- so
the speed advantage does not come at a cost in solution quality.

\begin{table}[ht]
  \centering
  \caption{Median wall-clock time (ms) over the five benchmark functions, by
  dimension. Lower is better.}
  \label{tab:time}
  \begin{tabular}{lrrr}
    \toprule
    Library & $d=2$ & $d=10$ & $d=30$ \\
    \midrule
    \texttt{turboswarm} (native) & \textbf{3.23} & \textbf{5.25} & \textbf{9.85} \\
    \texttt{turboswarm} (Python) & 43.87 & 50.44 & 58.48 \\
    \texttt{pyswarms}            & 10.43 & 13.40 & 20.52 \\
    \texttt{pyswarm}             & 45.61 & 64.81 & 67.41 \\
    DEAP                         & 46.01 & 64.09 & 106.55 \\
    \texttt{pymoo}               & 225.89 & 227.38 & 241.71 \\
    \bottomrule
  \end{tabular}
\end{table}

\begin{table}[ht]
  \centering
  \caption{Speedup of \texttt{turboswarm} (native path) over each baseline
  (median time ratio over the five functions). Higher is better.}
  \label{tab:speedup}
  \begin{tabular}{lrrr}
    \toprule
    Baseline & $d=2$ & $d=10$ & $d=30$ \\
    \midrule
    \texttt{turboswarm} (Python) & $13.6\times$ & $9.6\times$ & $5.9\times$ \\
    \texttt{pyswarms}            & $3.2\times$  & $2.6\times$ & $2.1\times$ \\
    \texttt{pyswarm}             & $14.1\times$ & $12.3\times$ & $6.8\times$ \\
    DEAP                         & $14.2\times$ & $12.2\times$ & $10.8\times$ \\
    \texttt{pymoo}               & $69.8\times$ & $43.3\times$ & $24.5\times$ \\
    \bottomrule
  \end{tabular}
\end{table}

\begin{figure}[ht]
  \centering
  \includegraphics[width=0.8\linewidth]{speedup.png}
  \caption{Speedup of \texttt{turboswarm} (native objective path) over
  \texttt{pyswarms} across dimensions, for five benchmark functions.}
  \label{fig:speedup}
\end{figure}

The benchmark was run with \texttt{turboswarm} 0.2.2, \texttt{pyswarms} 1.3.0,
\texttt{pyswarm} 1.0.0, \texttt{pymoo} 0.6.1.6, DEAP 1.4 and NumPy 2.4.6 under
Python 3.12, on an Apple M5 (10 cores, macOS 26.5). Absolute numbers are
machine-dependent; the full benchmark --- including the raw results and machine
provenance --- is reproducible with \texttt{benches/bench\_suite.py}.

\section{Availability and citation}

\texttt{turboswarm} is released under the MIT license. The Python package is on
PyPI as \texttt{turboswarm} (\texttt{pip install turboswarm}) and the Rust core
on crates.io as \texttt{turboswarm-core}. The source code, documentation and
issue tracker are at \url{https://github.com/turboswarm/turboswarm.github.io},
with a documentation portal at \url{https://turboswarm.github.io}. Each release
is archived on Zenodo under the concept DOI \texttt{10.5281/zenodo.20832446}. If
you use \texttt{turboswarm} in academic work, please cite this paper and the
archived software release.

\section{Conclusion}

\texttt{turboswarm} provides a compiled, reproducible and extensible PSO toolkit
that is as easy to call as a pure-Python package. By exposing many PSO variants,
topologies and search spaces --- including an interval-valued ``grey'' space ---
behind one reproducible interface, it lowers the cost of comparative algorithmic
studies and supports teaching through its built-in visualization. Its
trait-based core is also the foundation for ongoing methodological work on
optimization with grey (interval-valued) variables.

\section*{AI usage disclosure}

The \texttt{turboswarm} library itself --- the Rust core, the Python bindings and
the algorithms --- was designed and written by the author. Generative AI
(Anthropic's Claude, used through Claude Code) assisted with drafting this paper,
implementing the comparative benchmark suite, and editing parts of the
documentation. All AI-assisted text and code were reviewed, tested and validated
by the author, who takes full responsibility for the content.

\section*{Acknowledgements}

The author received no specific funding for this work, and acknowledges the
open-source Rust and PyO3 communities, whose tooling underpins the Rust--Python
boundary used here.

\bibliographystyle{plainnat}
\bibliography{refs}

\end{document}
